% Generated from locale/tr/content/first-order-logic/tableaux/provability-consistency.tex; do not hand-edit.


\iftag{FOL}
      {\olfileid[tr]{fol}{tab}{prv}}
      {\olfileid[tr]{pl}{tab}{prv}}

\olsection{\usetoken{S}{derivability} ve Tutarlılık}

Şimdi !!{derivability} bağıntısının çeşitli özelliklerini belirleyeceğiz.
Bunlar kendi başlarına da ilginçtir; ayrıca her biri tamlık teoreminin
ispatında rol oynayacaktır.

\begin{prop}\ollabel{prop:provability-contr}
  $\Gamma \Proves !A$ ve $\Gamma \cup \{!A\}$ tutarsızsa $\Gamma$
  tutarsızdır.
\end{prop}

\begin{proof}
Aşağıdaki iki kümenin kapalı !!{tableau}s ağaçları olacak biçimde sonlu
$\Gamma_0 = \{!B_1, \dots, !B_n\} \subseteq \Gamma$ ve
$\Gamma_1 =\{!C_1, \dots, !C_m\} \subseteq \Gamma$ kümeleri vardır:
  \begin{align*}
    \{\sFmla{\False}{!A}, &
    \sFmla{\True}{!B_1}, \dots, \sFmla{\True}{!B_n}\}, \\
    \{\sFmla{\True}{!A}, &
    \sFmla{\True}{!C_1}, \dots, \sFmla{\True}{!C_m}\}.
  \end{align*}
  $!A$ üzerinde \Cut{} kuralını kullanarak bunları $\Gamma_0
  \cup \Gamma_1$ kümesinin tutarsız olduğunu gösteren tek bir kapalı
  !!{tableau} hâlinde birleştirebiliriz. $\Gamma_0 \subseteq \Gamma$ ve
  $\Gamma_1 \subseteq \Gamma$ olduğundan $\Gamma_0 \cup
  \Gamma_1 \subseteq \Gamma$; dolayısıyla $\Gamma$ tutarsızdır.
\end{proof}

\begin{prop}
\ollabel{prop:prov-incons}
$\Gamma \Proves !A$ ancak ve ancak $\Gamma \cup \{\lnot !A\}$ tutarsızsa.
\end{prop}

\begin{proof}
Önce $\Gamma \Proves !A$ olduğunu, yani
\[
\{\sFmla{\False}{!A},
\sFmla{\True}{!B_1}, \dots, \sFmla{\True}{!B_n}\}
\]
için kapalı bir !!{tableau} bulunduğunu varsayın. $\TRule{\True}{\lnot}$
kuralı kullanılarak bu,
\[
\{\sFmla{\True}{\lnot !A},
\sFmla{\True}{!B_1}, \dots, \sFmla{\True}{!B_n}\}.
\]
için kapalı bir !!{tableau}'ya dönüştürülebilir.

Tersine, ikinci küme için kapalı bir !!{tableau} varsa, ilk
varsayım~$\sFmla{\True}{\lnot !A}$ ile ona
\TRule{\True}{\lnot} uygulanarak elde edilen bütün formülleri kaldırıp
$\sFmla{\False}{!A}$ varsayımını ekleyerek onu ilk kümenin kapalı bir
!!{tableau}'suna dönüştürebiliriz. Nitekim bir dal daha önce
$\sFmla{\True}{\lnot !A}$ ifadesine \TRule{\True}{\lnot} uygulanmasının
sonucu olan $\sFmla{\False}{!A}$ ifadesini içerdiği için kapalıysa, yeni
!!{tableau}'daki karşılık gelen dal da kapalıdır. Eski tableau'daki bir dal
$\sFmla{\True}{\lnot !A}$ varsayımının yanı sıra
$\sFmla{\False}{\lnot !A}$ ifadesini içerdiği için kapalıysa,
$\sFmla{\False}{\lnot !A}$ ifadesine $\TRule{\False}{\lnot}$ uygulayıp
$\sFmla{\True}{!A}$ elde ederek onu kapalı bir dala dönüştürebiliriz.
$\sFmla{\False}{!A}$ ifadesini varsayım olarak eklediğimizden bu dal
kapanır.
\end{proof}

\begin{prob}
$\Gamma \Proves \lnot !A$ olmasının ancak ve ancak
$\Gamma \cup \{!A\}$ kümesinin tutarsız olmasına denk olduğunu ispatlayın.
\end{prob}

\begin{prop}\ollabel{prop:explicit-inc}
  $\Gamma \Proves !A$ ve $\lnot !A \in \Gamma$ ise $\Gamma$ tutarsızdır.
\end{prop}

\begin{proof}
  $\Gamma \Proves !A$ ve $\lnot !A \in \Gamma$ olduğunu varsayın. Bu
  durumda $!B_1$, \dots, $!B_n \in \Gamma$ vardır ve
  \[
  \{\sFmla{\False}{!A}, \sFmla{\True}{!B_1}, \dots, \sFmla{\True}{!B_n}\}
  \]
  kümesinin kapalı bir tableau'su bulunur. \sFmla{\False}{!A} varsayımını
  \sFmla{\True}{\lnot !A} ile değiştirin ve
  \sFmla{\True}{\lnot !A} ifadesine \TRule{\True}{\lnot} uygulanmasının
  sonucunu varsayımların arkasına ekleyin. Her !!{sentence}, !!{tableau}
  içinde $1$.~satıra dayanılarak gerekçelendirilmişse eski !!{tableau}'daki
  bu gerekçe artık $n+2$.~satıra dayanır. Dolayısıyla eski !!{tableau}
  kapalıysa yenisi de kapalıdır. Tüm varsayımlar~$\Gamma$ içinde olduğundan
  yeni tableau, $\Gamma$'nın tutarsız olduğunu gösterir.
\end{proof}

\begin{prop}\ollabel{prop:provability-exhaustive}
  $\Gamma \cup \{!A\}$ ve $\Gamma \cup \{\lnot !A\}$ kümelerinin ikisi de
  tutarsızsa $\Gamma$ tutarsızdır.
\end{prop}

\begin{proof}
  $!B_1$, \dots, $!B_n \in \Gamma$ ve $!C_1$, \dots,
  $!C_m \in \Gamma$ var ve
  \begin{align*}
    \{\sFmla{\True}{!A}, &
    \sFmla{\True}{!B_1}, \dots, \sFmla{\True}{!B_n}\} \text{ ile}\\
    \{\sFmla{\True}{\lnot !A}, &
    \sFmla{\True}{!C_1}, \dots, \sFmla{\True}{!C_m}\}
  \end{align*}
  kümelerinin ikisinin de kapalı !!{tableau}s ağaçları varsa, varsayımlar
  olarak $\sFmla{\True}{!B_1}$, \dots, $\sFmla{\True}{!B_n}$ ile
  $\sFmla{\True}{!C_1}$, \dots, $\sFmla{\True}{!C_m}$ ifadelerini birlikte
  kullanıp ardından \Cut{} kuralını uygulayarak $\Gamma$'nın tutarsız
  olduğunu gösteren tek bir birleşik !!{tableau} kurabiliriz. Böylece biri
  $\sFmla{\True}{!A}$, diğeri $\sFmla{\False}{!A}$ ile başlayan iki dal
  elde edilir.
  
  Sol tarafa, ilk !!{tableau}'nun varsayımlarının altında kalan kısmını
  ekleyin. İlk !!{tableau}'nun $\sFmla{\True}{!A}$ dâhil her varsayımı
  bulunduğundan, buradaki bütün kural uygulamaları hâlâ doğrudur. Böylece
  $\sFmla{\True}{!A}$ altındaki her dal kapanır.
  
  Sağ tarafa, ikinci !!{tableau}'nun varsayımının altında kalan kısmını;
  $\sFmla{\True}{\lnot !A}$ ifadesine~$\TRule{\True}{\lnot}$ kuralının
  uygulanmasıyla elde edilen sonuçları kaldırarak ekleyin.
  $\sFmla{\True}{\lnot !A}$ ifadesine $\TRule{\True}{\lnot}$ uygulamanın
  sonucu $\sFmla{\False}{!A}$'dır; bu ifade birleşik !!{tableau}'nun sağ
  tarafında \Cut{} kuralının sonucu olarak zaten bulunmaktadır.

  İkinci tableau'daki bir dal, birleşik !!{tableau}'da artık varsayım olarak
  yer almayan $\sFmla{\True}{\lnot !A}$ varsayımının yanı sıra
  $\sFmla{\False}{\lnot !A}$ ifadesini içerdiği için kapanmışsa,
  $\sFmla{\False}{\lnot !A}$ ifadesine $\TRule{\False}{\lnot}$ uygulayıp
  $\sFmla{\True}{!A}$ elde edebiliriz. Birleşik !!{tableau}'daki karşılık
  gelen dal artık kapanır; çünkü \Cut{} kuralının sağdaki sonucu olan
  $\sFmla{\False}{!A}$'yı içerir. İkinci !!{tableau}'daki bir dal başka bir
  nedenle kapanmışsa birleşik !!{tableau}'daki karşılık gelen dal da
  kapanır; çünkü $\sFmla{\True}{\lnot !A}$ dışındaki her !!{signed formula}s,
  eski ikinci !!{tableau}'nun dalında geçtiği gibi birleşik !!{tableau}'nun
  karşılık gelen dalında da geçer.
  \end{proof}

